\chapter{PROPOSED SYSTEM/ METHODOLOGY}
\section{System Overview}

MotionBridge is a four-stage real-time pipeline. In the first stage, three USB webcams positioned at front, left, and right angles simultaneously capture video of a human subject performing arm movements. In the second stage, MediaPipe Pose and MediaPipe Hands run independently on each camera feed to extract 2D joint keypoints including shoulder, elbow, wrist, and 21 finger landmarks along with per-joint visibility confidence scores. In the third stage, the 2D detections from all three cameras are fused through confidence-weighted triangulation using calibrated projection matrices to produce 3D joint coordinates in metric world space. In the fourth stage, the reconstructed 3D wrist position is fed into PyBullet's Inverse Kinematics solver, which computes the corresponding joint angles for a simulated arm model and updates the simulation every frame. The resulting 3D coordinates and joint angles are simultaneously written to a CSV dataset file.

\section*{Block Diagram / Flowchart}

\begin{figure}[h!]
\centering
\begin{tikzpicture}[
    node distance=1.2cm,
    block/.style={rectangle, draw, text width=4.5cm, align=center, rounded corners, minimum height=1cm, font=\small},
    input/.style={rectangle, draw, text width=2.5cm, align=center, minimum height=0.8cm, font=\small},
    output/.style={rectangle, draw, text width=3.5cm, align=center, minimum height=1cm, font=\small},
    arrow/.style={thick,->,>=stealth}
]

% Input Nodes (Side-by-side at the top)
\node (cam2) [input] {Camera 2\\(Left)};
\node (cam1) [input, left=0.5cm of cam2] {Camera 1\\(Front)};
\node (cam3) [input, right=0.5cm of cam2] {Camera 3\\(Right)};

% Processing Nodes (Stacked vertically)
\node (mediapipe) [block, below=1.2cm of cam2] {MediaPipe\\Pose \& Hands};
\node (triangulation) [block, below=of mediapipe] {Confidence-Weighted Triangulation};
\node (pybullet) [block, below=of triangulation] {PyBullet IK Solver};

% Output Nodes (Side-by-side at the bottom)
\node (sim) [output, below left=1.2cm and 0.2cm of pybullet.south] {Real-Time Simulation};
\node (csv) [output, below right=1.2cm and 0.2cm of pybullet.south] {Structured CSV Dataset};

% Drawing Arrows from Inputs to MediaPipe
\draw [arrow] (cam1.south) -- ++(0,-0.4) -| (mediapipe.north);
\draw [arrow] (cam2.south) -- (mediapipe.north);
\draw [arrow] (cam3.south) -- ++(0,-0.4) -| (mediapipe.north);

% Drawing Arrows between Processing Blocks
\draw [arrow] (mediapipe.south) -- (triangulation.north);
\draw [arrow] (triangulation.south) -- (pybullet.north);

% Drawing Arrows from IK Solver to Outputs
\draw [arrow] (pybullet.south) -- ++(0,-0.4) -| (sim.north);
\draw [arrow] (pybullet.south) -- ++(0,-0.4) -| (csv.north);

\end{tikzpicture}
\caption{MotionBridge Real-Time Distributed Software Pipeline}
\label{fig:pipeline_vertical}
\end{figure}

\section{Algorithms and Working Principle}

\subsection{Camera Calibration}

Prior to operation, each camera is calibrated using OpenCV's checkerboard calibration procedure to obtain the intrinsic matrix $\mathbf{K}$ (focal lengths, principal point, distortion coefficients) and extrinsic parameters $[\mathbf{R} | \mathbf{t}]$ (rotation and translation relative to the world coordinate frame). The projection matrix for each camera $i$ is computed as:

\begin{equation}
    \mathbf{P}_i = \mathbf{K}_i \cdot [\mathbf{R}_i | \mathbf{t}_i]
\end{equation}

\subsection{MediaPipe Pose and Hand Estimation}

For each video frame from camera $i$, MediaPipe Pose returns a set of 33 body landmarks and MediaPipe Hands returns 21 hand landmarks, each with associated 2D pixel coordinates $(u, v)$ and a visibility confidence score $c \in [0, 1]$. The arm joints of interest are: shoulder, elbow, wrist, and all 21 finger keypoints.

\subsection{Confidence-Weighted Triangulation}

Standard linear triangulation computes the 3D position $\mathbf{X}$ of a joint from its 2D observations across $N$ calibrated cameras by solving a linear system $\mathbf{A}\mathbf{X} = 0$ via Singular Value Decomposition (SVD). The proposed system extends this with confidence weighting, scaling each camera's contribution by its MediaPipe visibility score:

\begin{equation}
    \mathbf{X}_{3D} = \frac{\sum_{i=1}^{N} c_i \cdot f(\mathbf{P}_i, \mathbf{x}_i)}{\sum_{i=1}^{N} c_i}
\end{equation}

where $c_i$ is the confidence score from camera $i$, $\mathbf{P}_i$ is its projection matrix, and $\mathbf{x}_i = (u_i, v_i)$ is the observed 2D joint location. This reduces the influence of cameras where a joint is occluded or detected with low confidence, improving reconstruction accuracy without additional hardware.

\subsection{Inverse Kinematics in PyBullet}

The reconstructed 3D wrist position $\mathbf{X}_{wrist}$ is passed as the target end-effector position to PyBullet's built-in Damped Least Squares IK solver:
\begin{equation}
    \boldsymbol{\theta} = \texttt{pybullet.calculateInverseKinematics}(\text{arm}, \text{endEffector}, \mathbf{X}_{wrist})
\end{equation}

The resulting joint angle vector $\boldsymbol{\theta}$ is applied to the simulated arm model each frame via \texttt{pybullet.setJointMotorControl2}, producing a live visual replication of the human arm movement.

\subsection{Dataset Generation}

Each processed frame produces one record written to a CSV file:

\begin{equation}
    \text{Record} = \{t,\ \mathbf{X}_{sh},\ \mathbf{X}_{el},\ \mathbf{X}_{wr},\ \mathbf{X}_{f_1},\ \ldots,\ \mathbf{X}_{f_{21}},\ \boldsymbol{\theta}\}
\end{equation}

where $t$ is the timestamp, $\mathbf{X}_{sh}$, $\mathbf{X}_{el}$, $\mathbf{X}_{wr}$ are the 3D shoulder, elbow, and wrist coordinates, $\mathbf{X}_{f_k}$ are the 21 finger joint coordinates, and $\boldsymbol{\theta}$ is the simulated joint angle vector.
